% Updated from the user-supplied Appendix A draft (attachment ff144394).
% Explicit appendix equation tags preserve A.1--A.32 without renumbering the main text.

This appendix establishes the results in Sections~\ref{sec:geometry}--\ref{sec:tail-identifiability}. We first derive the nuisance-projected information and its tangency construction. We then prove the inverse-information tail law, the local classification lower bound, and the information-coverage property of Fisher-Glass. Throughout, task coordinates use the normalization fixed in the main text, and environmental aggregation is with respect to the declared distribution \(\Pi\).

\subsection{Projection, Normalization, and Task--Nuisance Tangency}
\label{app:theory-projection}

Fix an environment \(e\) and its reference distribution \(P_e\). Let
\[
s_{T,e}\in L_0^2(P_e)^{d_T},
\qquad
s_{N,e}\in L_0^2(P_e)^{d_N}
\]
be centered, square-integrable task and nuisance scores. Write
\[
J_{AB,e}
=
\mathbb E_{P_e}\!\left[s_{A,e}s_{B,e}^{\top}\right],
\qquad A,B\in\{T,N\}.
\]
Assume that the coordinates of \(s_{N,e}\) span the specified nuisance tangent space \(\mathcal T_{N,e}\) and that \(J_{NN,e}\) is positive definite.

\paragraph{Efficient-score projection.}

Set \(\Gamma_e=J_{TN,e}J_{NN,e}^{-1}\). The residual satisfies
\[
\mathbb E_{P_e}\!\left[
(s_{T,e}-\Gamma_es_{N,e})s_{N,e}^{\top}
\right]
=
J_{TN,e}-\Gamma_eJ_{NN,e}
=
0.
\tag{A.1}\label{eq:app-score-orthogonality}
\]
Because \(\Gamma_es_{N,e}\) lies componentwise in \(\mathcal T_{N,e}\), this orthogonality identifies it as the componentwise \(L^2(P_e)\) projection of the task score. Consequently,
\[
\begin{aligned}
s_{T,e}^{\mathrm{eff}}
&=
s_{T,e}-J_{TN,e}J_{NN,e}^{-1}s_{N,e},
\\
F_e
&=
\mathbb E_{P_e}\!\left[
s_{T,e}^{\mathrm{eff}}
(s_{T,e}^{\mathrm{eff}})^{\top}
\right]
\\
&=
J_{TT,e}-J_{TN,e}J_{NN,e}^{-1}J_{NT,e}.
\end{aligned}
\tag{A.2}\label{eq:app-efficient-schur}
\]
This is the efficient-information construction used in Equation~\eqref{eq:efficient-score-information} \citep{bickel1993efficient,tsiatis2006semiparametric}. Equivalently, for every task direction \(u\),
\[
\begin{aligned}
u^{\top}F_eu
&=
\min_{v\in\mathbb R^{d_N}}
\mathbb E_{P_e}\!\left[
(u^{\top}s_{T,e}-v^{\top}s_{N,e})^2
\right],
\\
v^\star
&=
J_{NN,e}^{-1}J_{NT,e}u.
\end{aligned}
\tag{A.3}\label{eq:app-variational-schur}
\]
Thus \(F_e\succeq0\). A nonzero direction \(u\) has zero efficient information exactly when \(u^{\top}s_{T,e}\) belongs to the nuisance span in \(L^2(P_e)\). The quantity \(\lambda_e=\lambda_{\min}(F_e)\) measures the weakest remaining first-order task information. For an infinite-dimensional nuisance tangent space, the orthogonal-projection definition applies directly to its closed span.

\paragraph{Coordinate normalization.}

Consider separate invertible score transformations
\[
\widetilde s_{T,e}=Us_{T,e},
\qquad
\widetilde s_{N,e}=Vs_{N,e}.
\]
The nuisance span is unchanged, so
\[
\widetilde s_{T,e}^{\mathrm{eff}}
=
Us_{T,e}^{\mathrm{eff}},
\qquad
\widetilde F_e
=
UF_eU^{\top}.
\tag{A.4}\label{eq:app-coordinate-congruence}
\]
An invertible nuisance-basis change therefore preserves exact task information when the task coordinates remain fixed. Task rescaling transforms information by congruence, preserving rank while generally changing its numerical eigenvalues. Comparisons of \(\lambda_e\) consequently use a common task-coordinate convention.

For one task and one nuisance score with positive raw information, define
\[
\rho_e
=
\frac{J_{TN,e}}{\sqrt{J_{TT,e}J_{NN,e}}}.
\]
Then
\[
F_e
=
J_{TT,e}(1-\rho_e^2),
\qquad
\frac{F_e}{J_{TT,e}}
=
1-\rho_e^2
\in[0,1].
\tag{A.5}\label{eq:app-surviving-fraction}
\]
Nonzero scalar rescaling of either score leaves this fraction unchanged. More generally, for a scalar task, the fraction is invariant under an invertible nuisance-basis change. These identities concern the exact projection; Appendix~\ref{app:estimation-certification} specifies the coordinate conventions for its ridge implementation.

The following construction expresses the transverse-information geometry of \citet{mansouri_arai_2026_fisher_glasses} in the normalized score coordinates used here.

\begin{appendixproposition}[Task--nuisance tangency]
\label{prop:task-nuisance-tangency}

Let \(b_e,r_e\) be centered, orthogonal, unit-norm elements of \(L^2(P_e)\). Suppose
\[
\mathcal T_{N,e}
=
\operatorname{span}\{b_e\}
\]
and the scalar task score is
\[
a_e=b_e+\delta_er_e,
\qquad
\delta_e\ge0.
\]
Then
\[
J_{TT,e}=1+\delta_e^2,
\qquad
F_e=\delta_e^2,
\qquad
\frac{F_e}{J_{TT,e}}
=
\frac{\delta_e^2}{1+\delta_e^2}.
\tag{A.6}\label{eq:app-tangency-identities}
\]

\end{appendixproposition}
\begin{proof} The coefficient of the projection of \(a_e\) onto \(b_e\) is
\[
\frac{\langle a_e,b_e\rangle_e}{\|b_e\|_e^2}=1.
\]
Hence its efficient score is \(\delta_er_e\), whose squared norm is \(\delta_e^2\). Orthogonality gives \(\|a_e\|_e^2=1+\delta_e^2\), and division yields the surviving fraction. \end{proof}

The construction keeps raw task information at least one while projected information vanishes quadratically. When the nuisance feature supports a predictor with small observed risk, that predictive success coexists with the displayed information collapse. This is the separation used in Section~\ref{sec:tangency-separation}; Appendix~\ref{app:synthetic-checks} specifies its synthetic realization.

\subsection{Normal Coordinates and the Inverse-Information Tail}
\label{app:theory-tails}

We prove the geometric tail law of \citet{mansouri_arai_2026_fisher_glasses} under Theorem~\ref{thm:inverse-information-geometry}'s assumptions, including its asymptotic constant. Let
\[
\lambda:\mathbb R^d\longrightarrow[0,\infty),
\]
and let
\[
\mathcal Z=\{z:\lambda(z)=0\}
\]
be a nonempty compact \(C^2\) embedded submanifold without boundary, of codimension \(c\in\{1,\ldots,d\}\). Choose a tubular radius \(r_\star>0\) and write
\[
U_{r_\star}(\mathcal Z)
=
\{z:\operatorname{dist}(z,\mathcal Z)<r_\star\}.
\]
Assume that \(\lambda\) is continuous on \(\overline{U_{r_\star}(\mathcal Z)}\) and bounded away from zero outside \(U_{r_\star}(\mathcal Z)\).

Let \(p\) be the environment density with respect to Lebesgue measure. For \(s\in\mathcal Z\), let \(S(N_s\mathcal Z)\) denote the unit sphere in its normal space. In normal coordinates \(z=s+r\omega\), assume uniformly on the unit normal bundle that
\[
\begin{aligned}
\lambda(s+r\omega)
&=
a(s,\omega)r^{2m}(1+o(1)),
\\
p(s+r\omega)
&=
b(s,\omega)r^\nu(1+o(1)),
\end{aligned}
\qquad r\downarrow0,
\tag{A.7}\label{eq:null-set-expansions}
\]
where \(m>0\), \(\nu>-c\), and \(a,b\) are continuous and bounded above and below by positive constants. Define
\[
\begin{aligned}
\beta
&=
\frac{c+\nu}{2m},
\\
C_{\mathcal Z}
&=
\frac{1}{c+\nu}
\int_{\mathcal Z}
\int_{S(N_s\mathcal Z)}
b(s,\omega)a(s,\omega)^{-\beta}
\,d\sigma_s(\omega)\,
d\operatorname{vol}_{\mathcal Z}(s),
\end{aligned}
\tag{A.8}\label{eq:app-tail-constant}
\]
where \(\sigma_s\) is the usual, unnormalized surface measure on the normal unit sphere and \(\operatorname{vol}_{\mathcal Z}\) is induced volume on \(\mathcal Z\). In codimension one, the sphere consists of the two normal directions, each with unit mass.

\begin{proof}[Proof of Theorem~\ref{thm:inverse-information-geometry}]

First localize the event of small information. For any \(r_0\in(0,r_\star)\), the annulus
\[
\left\{
z:
r_0\le\operatorname{dist}(z,\mathcal Z)\le r_\star
\right\}
\]
is compact and contains no zeros of \(\lambda\). Continuity gives a positive minimum there. Together with the lower bound outside the outer tube, this implies that
\[
\{\lambda<t\}\subset U_{r_0}(\mathcal Z)
\]
for all sufficiently small \(t>0\).

Within the tube, the volume element is
\[
\begin{aligned}
dz
&=
j(s,r\omega)\,r^{c-1}
\,dr\,d\sigma_s(\omega)\,
d\operatorname{vol}_{\mathcal Z}(s),
\\
j(s,r\omega)
&=
1+O(r),
\end{aligned}
\tag{A.9}\label{eq:app-tube-jacobian}
\]
uniformly on the unit normal bundle. Fix \(\zeta\in(0,1)\). The uniform expansions and the Jacobian estimate allow \(r_0\) to be chosen so that
\[
\begin{aligned}
(1-\zeta)a(s,\omega)r^{2m}
&\le
\lambda(s+r\omega)
\le
(1+\zeta)a(s,\omega)r^{2m},
\\
(1-\zeta)b(s,\omega)r^\nu
&\le
p(s+r\omega)j(s,r\omega)
\le
(1+\zeta)b(s,\omega)r^\nu
\end{aligned}
\tag{A.10}\label{eq:app-uniform-tube-bounds}
\]
for \(0<r<r_0\).

For sufficiently small \(t\), all relevant radial thresholds lie within this tube. The event \(\lambda<t\) contains the radial region
\[
r<
\left(
\frac{t}{(1+\zeta)a(s,\omega)}
\right)^{1/(2m)}
\]
and is contained in
\[
r<
\left(
\frac{t}{(1-\zeta)a(s,\omega)}
\right)^{1/(2m)}.
\]
Integrating the corresponding bounds gives
\[
\begin{aligned}
(1-\zeta)(1+\zeta)^{-\beta}C_{\mathcal Z}t^\beta
&\le
\Pr(\lambda<t)
\\
&\le
(1+\zeta)(1-\zeta)^{-\beta}C_{\mathcal Z}t^\beta.
\end{aligned}
\tag{A.11}\label{eq:app-small-information-sandwich}
\]
Here the radial integral is
\[
\int_0^{(t/a)^{1/(2m)}}
b\,r^{c+\nu-1}\,dr
=
\frac{b\,a^{-\beta}}{c+\nu}\,t^\beta.
\]
Letting \(t\downarrow0\) and then \(\zeta\downarrow0\) proves
\[
\Pr(\lambda<t)\sim C_{\mathcal Z}t^\beta.
\]
Compactness and the positive bounds on \(a,b\) give \(0<C_{\mathcal Z}<\infty\). Since \(c\ge1\), the null manifold has zero Lebesgue measure and hence zero environment probability. Taking \(t=x^{-1}\) therefore yields
\[
\Pr\!\left(\lambda(z)^{-1}>x\right)
\sim
C_{\mathcal Z}x^{-\beta}.
\tag{A.12}\label{eq:app-geometric-tail}
\]
This proves Theorem~\ref{thm:inverse-information-geometry}. \end{proof}

\paragraph{Inverse-information moments.}

For every \(p_0>0\), the tail-integral identity gives
\[
\begin{aligned}
\mathbb E[\lambda^{-p_0}]
&=
p_0\int_0^\infty
x^{p_0-1}\Pr(\lambda^{-1}>x)\,dx,
\\
\mathbb E[\lambda^{-p_0}]<\infty
&\quad\Longleftrightarrow\quad
p_0<\beta.
\end{aligned}
\tag{A.13}\label{eq:app-inverse-moments}
\]
The integral over bounded \(x\) is finite, and its tail has the convergence criterion of
\[
\int^\infty x^{p_0-\beta-1}\,dx.
\]
In particular, the exact inverse-information mean is finite precisely when \(\beta>1\); at \(\beta=1\), the divergence is logarithmic. Appendix~\ref{app:theory-certificate} gives the corresponding certificate limit.

\paragraph{Tangency and independent information branches.}

For the scalar construction \(\lambda=\delta^2\),
\[
\Pr(\delta<t)\sim C_\delta t^\gamma
\]
immediately gives
\[
\Pr(\lambda^{-1}>x)
\sim
C_\delta x^{-\gamma/2}.
\]

The nonnegative residual amplitudes used in the synthetic experiments admit the following scalar portfolio specialization of \citet{mansouri_arai_2026_fisher_glasses}. Suppose
\[
\lambda_K=\sum_{k=1}^K w_k\delta_k^2,
\qquad w_k>0,
\]
and \(\delta\in[0,\infty)^K\) has a joint density \(f\) continuous and positive at the origin from within the orthant. A change of variables gives
\[
\begin{aligned}
\Pr(\lambda_K<t)
&=
\frac{t^{K/2}}{\sqrt{\prod_{k=1}^K w_k}}
\int_{B_K^+}
f\!\left(
\sqrt{t/w_1}\,u_1,\ldots,
\sqrt{t/w_K}\,u_K
\right)\,du
\\
&\sim
\frac{f(0)v_K}
{2^K\sqrt{\prod_{k=1}^K w_k}}\,
t^{K/2},
\end{aligned}
\tag{A.14}\label{eq:app-independent-branches}
\]
where \(B_K^+\) is the positive-orthant portion of the unit ball and \(v_K\) is the volume of the full unit ball. Independence gives
\[
f(0)=\prod_{k=1}^K f_k(0+).
\]
The inverse-information exponent is therefore \(K/2\), and its first moment is finite for \(K\ge3\).

This calculation applies to the specified additive scalar information, or to information in a fixed task direction. A minimum-eigenvalue conclusion is governed by Theorem~\ref{thm:inverse-information-geometry} applied to that eigenvalue. Appendix~\ref{app:synthetic-checks} evaluates the scalar construction.

\subsection{Local Classification: Testing Reduction and Tail Aggregation}
\label{app:theory-classification}

For a distribution \(P\), let
\[
\mathcal E_P(g)=R_P(g)-R_P^\star
\]
denote population excess classification risk. Fix an environment \(e\) and a local family \(\mathcal P_e^{\mathrm{loc}}\), specified independently of the sample size. Its minimax expected excess risk is
\[
\mathfrak R_e(n)
=
\inf_{\widehat g_n}
\sup_{P\in\mathcal P_e^{\mathrm{loc}}}
\mathbb E_{S_n\sim P^{\otimes n}}
\!\left[
\mathcal E_P\bigl(\widehat g_n(S_n)\bigr)
\right],
\tag{A.15}\label{eq:app-local-minimax}
\]
where the infimum ranges over measurable learning rules. The sample consists of \(n\) independent observations from environment \(e\). For randomized rules, the expectation also includes their independent randomization.

Assume \(\lambda_e>0\). For every \(0<t\le t_{0,e}\), the local family contains \(P_{e,+t}\) and \(P_{e,-t}\) with task coordinates \(\pm tu_e\) along a unit weakest efficient-information direction; nuisance coordinates may vary with the alternative. For constants \(K_e,c_e,t_{0,e}>0\) and \(q>0\), suppose
\[
\operatorname{KL}(P_{e,+t}\Vert P_{e,-t})
\le K_et^2\lambda_e
\]
and
\[
\inf_g
\left\{
\mathcal E_{P_{e,+t}}(g)
+
\mathcal E_{P_{e,-t}}(g)
\right\}
\ge c_et^q
\]
throughout this radius, as in Theorem~\ref{thm:local-classification-difficulty}.

\begin{proof}[Proof of Theorem~\ref{thm:local-classification-difficulty}]

Fix \(0<t\le t_{0,e}\) and abbreviate \(P_+=P_{e,+t}\) and \(P_-=P_{e,-t}\). The decision-separation assumption gives
\[
\mathcal E_{P_+}(g)+\mathcal E_{P_-}(g)\ge c_et^q
\]
for every classifier \(g\). For an arbitrary learning rule, define the test
\[
\psi_t(S_n)
=
\begin{cases}
+,
&
\mathcal E_{P_+}(\widehat g_n(S_n))
\le
\mathcal E_{P_-}(\widehat g_n(S_n)),
\\
-,
&
\text{otherwise}.
\end{cases}
\tag{A.16}\label{eq:app-induced-test}
\]
If this test selects the wrong alternative, the excess risk under the true alternative is at least \(c_et^q/2\). Thus
\[
\begin{aligned}
&\max_{\sigma\in\{+,-\}}
\mathbb E_{P_\sigma^{\otimes n}}
\!\left[
\mathcal E_{P_\sigma}(\widehat g_n)
\right]
\\
&\qquad\ge
\frac{c_et^q}{4}
\left\{
P_+^{\otimes n}(\psi_t=-)
+
P_-^{\otimes n}(\psi_t=+)
\right\}.
\end{aligned}
\tag{A.17}\label{eq:app-risk-to-testing}
\]
The standard two-point testing inequality gives
\[
\begin{aligned}
&P_+^{\otimes n}(\psi_t=-)
+
P_-^{\otimes n}(\psi_t=+)
\\
&\qquad\ge
\frac12
\exp\!\left[
-\operatorname{KL}
\!\left(
P_+^{\otimes n}\Vert P_-^{\otimes n}
\right)
\right]
\end{aligned}
\tag{A.18}\label{eq:app-testing-errors}
\]
\citep{lecam1986asymptotic,tsybakov2009introduction}. This inequality follows, for example, from the Hellinger affinity \(a(P,Q)=\int\sqrt{pq}\): Jensen's inequality gives
\[
a(P,Q)\ge
\exp[-\operatorname{KL}(P\Vert Q)/2],
\]
Cauchy--Schwarz gives
\[
\operatorname{TV}(P,Q)\le\sqrt{1-a(P,Q)^2},
\]
and the sum of testing errors is at least
\[
1-\operatorname{TV}(P,Q)\ge a(P,Q)^2/2.
\]

Independence and the assumed KL bound imply
\[
\operatorname{KL}
\!\left(P_+^{\otimes n}\Vert P_-^{\otimes n}\right)
=
n\operatorname{KL}(P_+\Vert P_-)
\le
nK_et^2\lambda_e.
\]
Both alternatives belong to the full local family. Taking its supremum and then the infimum over learning rules yields, for every admissible \(t\),
\[
\mathfrak R_e(n)
\ge
\frac{c_e}{8}t^q
\exp(-nK_et^2\lambda_e).
\tag{A.19}\label{eq:app-testing-reduction}
\]
Choose
\[
t_n
=
\min\left\{
t_{0,e},
(nK_e\lambda_e)^{-1/2}
\right\}.
\]
Then \(nK_e\lambda_et_n^2\le1\), and hence
\[
\mathfrak R_e(n)
\ge
c_0c_e
\min\left\{
t_{0,e}^{q},
(nK_e\lambda_e)^{-q/2}
\right\},
\qquad
c_0=\frac{\exp(-1)}8.
\tag{A.20}\label{eq:app-local-risk-bound}
\]
This proves Theorem~\ref{thm:local-classification-difficulty}. Independent learner randomization can be included in the testing experiment with the same law under both alternatives, leaving the KL calculation unchanged. \end{proof}

\paragraph{Quadratic decision separation.}

A sufficient condition for \(q=2\) can be expressed through decision-boundary mass. Suppose the binary alternatives share a feature marginal \(Q_e\), with regression functions
\[
\eta_{\pm,t}(x)
=
P_{e,\pm t}(Y=1\mid X=x)
\]
and Bayes rules \(g_{\pm,t}^\star\). On their disagreement set
\[
D_{e,t}
=
\{x:g_{+,t}^\star(x)\ne g_{-,t}^\star(x)\},
\]
any classifier disagrees with at least one Bayes rule. The binary excess-risk identity gives
\[
\begin{aligned}
&\inf_g
\left\{
\mathcal E_{P_{e,+t}}(g)
+
\mathcal E_{P_{e,-t}}(g)
\right\}
\\
&\qquad=
\int_{D_{e,t}}
\min\left\{
|2\eta_{+,t}(x)-1|,
|2\eta_{-,t}(x)-1|
\right\}\,dQ_e(x).
\end{aligned}
\tag{A.21}\label{eq:app-boundary-separation}
\]
Equality holds for unrestricted measurable binary classifiers; the integral is a lower bound for any restricted class.

If, for every \(0<t\le t_{0,e}\), a measurable \(A_{e,t}\subseteq D_{e,t}\) satisfies
\[
Q_e(A_{e,t})\ge c_{\mathrm{mass},e}t
\]
and both margins are at least \(c_{\mathrm{margin},e}t\) on \(A_{e,t}\), for positive constants independent of \(t\), then the right-hand side is at least
\[
c_{\mathrm{mass},e}c_{\mathrm{margin},e}t^2.
\]
Thus \(q=2\) and
\[
c_e=c_{\mathrm{mass},e}c_{\mathrm{margin},e}
\]
are valid. This is the moving-boundary regime in which disagreement mass and the relevant margins each scale linearly with the perturbation. Combined with the KL condition of Theorem~\ref{thm:local-classification-difficulty}, it gives the inverse-information factor in the quadratic local bound.

\paragraph{Aggregation across environments.}

Suppose the assumptions hold for \(\Pi\)-almost every environment and the resulting quantities are measurable. Define
\[
\mathcal B_e(n)
=
c_e
\min\left\{
t_{0,e}^{q},
(nK_e\lambda_e)^{-q/2}
\right\}.
\tag{A.22}\label{eq:app-local-bound-factor}
\]
Monotonicity and positive homogeneity of CVaR give
\[
\operatorname{CVaR}_{\alpha,e\sim\Pi}
\!\left(\mathfrak R_e(n)\right)
\ge
c_0\,
\operatorname{CVaR}_{\alpha,e\sim\Pi}
\!\left(\mathcal B_e(n)\right).
\tag{A.23}\label{eq:app-tail-risk-aggregation}
\]
The environment-specific separation strength, distinguishability constant, and local radius remain inside the aggregation. For \(q=2\), the interior regime
\[
nK_e\lambda_e\ge t_{0,e}^{-2}
\]
gives \(\mathcal B_e(n)=c_e/(nK_e\lambda_e)\); the local radius controls the bound beyond that regime. These are the statistical factors retained in Section~\ref{sec:fg-certificate} when the inverse-information tail motivates the audit.

\subsection{CVaR, Information Coverage, and Stabilization}
\label{app:theory-certificate}

We use the Fisher-glass certificate of \citet{mansouri_arai_2026_fisher_glasses}, defined through upper-tail CVaR. Fix \(0<\alpha<1\). For a nonnegative random variable \(Z\), allowing extended values, upper-tail conditional value-at-risk is
\[
\operatorname{CVaR}_{\alpha}(Z)
=
\inf_{\tau\in\mathbb R}
\left\{
\tau+
\frac{1}{1-\alpha}
\mathbb E[(Z-\tau)_+]
\right\}
\tag{A.24}\label{eq:cvar-definition}
\]
\citep{rockafellar2000optimization,rockafellar2002conditional}. Its upper-tail weighting representation is
\[
\begin{aligned}
\operatorname{CVaR}_{\alpha}(Z)
&=
\sup_{\xi\in\mathcal W_\alpha}
\mathbb E[\xi Z],
\\
\mathcal W_\alpha
&=
\left\{
\xi:
0\le\xi\le\frac1{1-\alpha},
\ \mathbb E\xi=1
\right\}.
\end{aligned}
\tag{A.25}\label{eq:app-cvar-dual}
\]
Indeed, for every feasible \(\xi\) and every \(\tau\),
\[
\mathbb E[\xi Z]
=
\tau+\mathbb E[\xi(Z-\tau)]
\le
\tau+\frac{\mathbb E[(Z-\tau)_+]}{1-\alpha}.
\]
Let
\[
v_\alpha
=
\inf\{v:\Pr(Z\le v)\ge\alpha\}.
\]
For integrable \(Z\), equality is attained by assigning weight \((1-\alpha)^{-1}\) above \(v_\alpha\), zero below it, and the remaining weight at \(v_\alpha\). Such a choice exists because
\[
\Pr(Z>v_\alpha)
\le
1-\alpha
\le
\Pr(Z\ge v_\alpha).
\]
Thus atoms at the threshold receive exactly the fraction needed to fill the upper tail. For nonnegative \(Z\) with infinite mean, both representations have value \(+\infty\).

This representation establishes monotonicity, positive homogeneity, and
\[
v_\alpha\le\operatorname{CVaR}_{\alpha}(Z).
\]
For a finite audit with weights \(\pi_g\ge0\) summing to one, it becomes
\[
\operatorname{CVaR}_{\alpha,\pi}(Z_g)
=
\max_{\substack{
\sum_g\omega_g=1\\
0\le\omega_g\le\pi_g/(1-\alpha)
}}
\sum_g\omega_g Z_g.
\tag{A.26}\label{eq:app-finite-cvar}
\]
These normalized upper-tail weights also provide the convention used in the information-sensitivity calculation of Appendix~\ref{app:tail-transverse-influence}.

\begin{proof}[Proof of Proposition~\ref{prop:information-coverage}]

Assume \(\varepsilon>0\) and \(\lambda_e\in[0,\infty)\) is finite almost surely. Write
\[
\begin{aligned}
L_e^{(\varepsilon)}
&=
(\lambda_e+\varepsilon)^{-1},
\\
I_\alpha^{\mathrm{FG}}(\varepsilon)
&=
\left[
\operatorname{CVaR}_{\alpha,e\sim\Pi}
\!\left(L_e^{(\varepsilon)}\right)
\right]^{-1}.
\end{aligned}
\tag{A.27}\label{eq:app-stabilized-certificate}
\]
The explicit argument \(\varepsilon\) records the stabilization used in Equation~\eqref{eq:fisher-glass-certificate}. If
\[
I_\alpha^{\mathrm{FG}}(\varepsilon)
\ge\rho>\varepsilon,
\]
then
\[
v_\alpha(L^{(\varepsilon)})
\le
\operatorname{CVaR}_{\alpha}(L^{(\varepsilon)})
\le
\rho^{-1}.
\]
By the quantile definition,
\[
\begin{aligned}
\Pr_{e\sim\Pi}\!\left(\lambda_e\ge\rho-\varepsilon\right)
&=
\Pr_{e\sim\Pi}\!\left(
L_e^{(\varepsilon)}\le\rho^{-1}
\right)
\\
&\ge\alpha.
\end{aligned}
\tag{A.28}\label{eq:app-information-coverage}
\]
This proves Proposition~\ref{prop:information-coverage}, including environments whose information equals the threshold. \end{proof}

Because \(0<L_e^{(\varepsilon)}\le\varepsilon^{-1}\), the stabilized certificate satisfies
\[
I_\alpha^{\mathrm{FG}}(\varepsilon)\ge\varepsilon.
\]
A level \(\rho>\varepsilon\) therefore certifies a strictly positive information floor \(\rho-\varepsilon\).

\paragraph{Finite-group interpretation.}

For four equally weighted groups and \(\alpha=0.95\), the group with the largest inverse-information value contains enough probability mass to fill the upper tail. Equation~\eqref{eq:app-finite-cvar} gives
\[
\begin{aligned}
\operatorname{CVaR}_{0.95}(L^{(\varepsilon)})
&=
\max_g L_g^{(\varepsilon)},
\\
I_{0.95}^{\mathrm{FG}}(\varepsilon)
&=
\min_g\lambda_g+\varepsilon.
\end{aligned}
\tag{A.29}\label{eq:app-four-group-certificate}
\]
The certificate identifies the weakest absolute information under the declared normalization. The surviving fraction \(F_g/J_{TT,g}\) measures relative task--nuisance alignment within a scalar-task group. The two quantities can therefore rank groups differently, as in Section~\ref{sec:prospective-diagnosis}.

\paragraph{The exact-information limit.}

Define \(L_e^{(0)}=\lambda_e^{-1}\), with \(1/0=+\infty\). As \(\varepsilon\downarrow0\),
\[
L_e^{(\varepsilon)}\uparrow L_e^{(0)}.
\]
For each feasible \(\xi\), monotone convergence gives
\[
\mathbb E[\xi L^{(\varepsilon)}]
\uparrow
\mathbb E[\xi L^{(0)}].
\]
Taking the supremum over \(\xi\) in Equation~\eqref{eq:app-cvar-dual} therefore yields
\[
\begin{aligned}
\operatorname{CVaR}_{\alpha}(L^{(\varepsilon)})
&\uparrow
\operatorname{CVaR}_{\alpha}(L^{(0)}),
\\
I_\alpha^{\mathrm{FG}}(\varepsilon)
&\downarrow
I_\alpha^{\mathrm{FG}}(0)
:=
\left[
\operatorname{CVaR}_{\alpha}(\lambda^{-1})
\right]^{-1},
\end{aligned}
\tag{A.30}\label{eq:app-exact-limit}
\]
where \(1/(+\infty)=0\). The interchange follows because the supremum over \(\xi\) and the increasing supremum along any sequence \(\varepsilon_k\downarrow0\) commute.

For any nonnegative \(Z\),
\[
\mathbb EZ
\le
\operatorname{CVaR}_{\alpha}(Z)
\le
\frac{\mathbb EZ}{1-\alpha}.
\tag{A.31}\label{eq:app-cvar-mean-bounds}
\]
The lower bound uses the feasible weight \(\xi=1\); the upper bound follows from \(\tau=0\) in the variational representation. Hence the exact certificate is strictly positive precisely when
\[
\mathbb E[\lambda^{-1}]<\infty.
\]
Under Theorem~\ref{thm:inverse-information-geometry},
\[
\begin{aligned}
I_\alpha^{\mathrm{FG}}(0)>0
&\quad\Longleftrightarrow\quad
\beta=\frac{c+\nu}{2m}>1,
\\
I_\alpha^{\mathrm{FG}}(0)=0
&\quad\text{for }\beta\le1.
\end{aligned}
\tag{A.32}\label{eq:app-certificate-integrability}
\]
This establishes the limiting interpretation in Section~\ref{sec:fg-certificate}. Information coverage is a population statement under \(\Pi\); the confidence level of an estimated lower certificate is supplied separately by the estimation analysis. Together, the geometric tail law, the conditional local-risk bound, and the coverage property define the population targets of the operational audit in Section~\ref{sec:computing-auditing}.
